\documentclass[10pt]{article}
\usepackage[margin=0.9in]{geometry}
\usepackage{amsmath,amssymb,amsthm}
\usepackage{booktabs,tabularx,array}
\usepackage{enumitem}
\usepackage{xcolor}
\usepackage[hidelinks]{hyperref}
\setlist{itemsep=2pt,topsep=3pt}

% Standalone citation macros: print the bib key so this file compiles without references.bib.
% Delete these three lines when moving text into the paper (natbib provides the real ones).
\providecommand{\citet}[1]{\textsf{#1}}
\providecommand{\citep}[1]{(\textsf{#1})}
\providecommand{\citealp}[1]{\textsf{#1}}

\newcommand{\eps}{\varepsilon}
\newcommand{\R}{\mathbb{R}}
\newcommand{\ind}{\mathbb{1}}
\newcommand{\ghat}{\hat{\mathbf g}}
\newcommand{\code}[1]{\texttt{\detokenize{#1}}}
\newtheorem{proposition}{Proposition}


\title{PhysLaw: condensed Related Work and Methods}
\date{}
\begin{document}
\sloppy
\maketitle

%=====================================================================
\section{Related Work (condensed)}
%=====================================================================

\paragraph{Violation of expectation and outcome prediction.}
Violation-of-expectation (VoE) benchmarks, borrowed from developmental psychology, test whether a model is more
surprised by an implausible video than by a matched plausible one. The violations are of qualitative object concepts:
objects disappearing (object permanence), changing shape (shape constancy, immutability), teleporting
(spatio-temporal continuity), passing through each other (solidity), or floating unsupported (gravity, support)
\citep{piloto2018probing,riochet2022intphys,bordes2025intphys2,weihs2022inflevel,jassim2024grasp}.
Physion \citep{bear2021physion} and Physion++ \citep{tung2023physionpp} instead train a readout to predict a binary
outcome (whether two cued objects make contact), where in Physion++ the outcome depends on latent mass, friction,
elasticity or deformability. \citet{garrido2025intuitive} find V-JEPA reliably above chance on such benchmarks while
pixel predictors and multimodal LLMs stay near chance. In every case the test is a yes/no judgement about whether an
event is possible; none asks whether a quantitative law holds.

\paragraph{Physical plausibility of generated videos.}
Output benchmarks score generated videos in three ways. Some ask a human or a vision-language model to rate physical
commonsense (VideoPhy \citep{bansal2024videophy}, PhyGenBench \citep{meng2024phygenbench}, VBench-2.0
\citep{zheng2025vbench2} and others). Some compare the generation with a ground-truth
continuation of the same conditioning frames: Physics-IQ \citep{motamed2025physicsiq} measures the overlap of motion
masks with real footage (where, when and how much motion occurs) plus pixel MSE, and RigidBench
\citep{jain2026rigidbench} measures the 3D trajectory error against a simulator. Others track objects in the generation
and check a physical quantity: PISA \citep{li2025pisa} measures the trajectory error of a falling object against real
drops, Morpheus \citep{tragoudaras2026morpheus} fits physics-informed models to the tracks and checks that conserved
quantities stay constant, and PVBench \citep{chen2026pvbench} scores generations with object-level physical measurements.
Measuring an absolute quantity (e.g.\ $g$ in m/s$^2$) from a generated video requires knowing its frame rate and
metric scale, which generations do not reliably have. The closest work therefore uses \emph{unit-free relational
tests}: two measurements of the same kind are taken in one video and compared as a ratio, so the unknown scale and
frame rate cancel. \citet{thozhiyoor2025slower} drop two objects from heights $h_1,h_2$ and test
$t_1^2/t_2^2=h_1/h_2$, where heights are in pixels and fall times in frames; Principia \citep{thozhiyoor2026principia}
generalises this to pairs of objects across eight phenomena, and PhysWeep \citep{khanbayov2026physweep} recovers the
physical parameter a generation realises. All of these evaluate only the generated output. None tests whether the
model's internal representations follow the law, so output-level and representation-level evidence for the same law
are never compared.

\paragraph{Physical variables inside video models.}
For encoders, \citet{joseph2026interpreting} find a ``Physics Emergence Zone'' at intermediate depth in V-JEPA~2 and
VideoMAE-v2, \citet{alam2026steering} steer a probe's possible/impossible classification of IntPhys videos, and
CALIPER \citep{mehta2026caliper} shows that in clean scenes even raw pixels predict how far a struck object slides. For diffusion transformers (DiTs), \citet{kneifl2026activations} decode
and steer speed, momentum and gravity, \citet{esmati2026invisible} decode plausibility, and LikePhys
\citep{yuan2026likephys} uses the denoising loss as a surprise score. This literature shows that individual physical variables (an object's speed, position or
plausibility label) can be decoded and sometimes steered. It does not test whether the relations among them obey a
law, for example whether decoded velocities before and after a collision conserve momentum. Outside video,
\citet{vafa2025inductive} show that a transformer that accurately predicts planetary orbits, when fine-tuned to predict
forces, implies a different and non-Newtonian law of gravitation on each slice of data: accurate prediction does not
imply that a model has learned the law behind it.

\paragraph{Physics post-training.}
PISA, VideoREPA \citep{zhang2025videorepa}, PhysRVG \citep{zhang2026physrvg}, NewtonRewards \citep{le2025newtonrewards}
and WROP \citep{zhang2026wrop} improve physics in generators and evaluate only outputs, so it is unknown whether
their output gains come with any change in what the model represents.

\paragraph{Interpretability benchmarks in other domains.}
Work on language and vision models has developed controls for claims about what a model encodes. Because probe
accuracy alone overstates what is represented \citep{belinkov2022probing}, probes are compared against control tasks
that keep the data's structure but remove the property (\emph{selectivity}, \citealp{hewitt2019control}) and against
the same architecture with random weights \citep{adebayo2018sanity}. Causal claims use interchange interventions,
which swap a representation between two inputs and check that the output changes as the hypothesised variable
predicts \citep{geiger2021causal}; RAVEL \citep{huang2024ravel}, CausalGym \citep{arora2024causalgym}
and MIB \citep{mueller2025mib} score localisation methods this way. Steering vectors are
known to vary in effect across inputs and sometimes push the wrong way \citep{tan2024steering,braun2025unreliability},
and steering along a fitted activation manifold behaves more naturally than linear steering
\citep{wurgaft2026manifold}. These benchmarks use language models, where the ground truth is a known circuit or a
labelled concept; none applies these controls to physical laws in video models.

%=====================================================================
\section{Methods}
\label{sec:methods}

\subsection{Overview}
PhysLaw asks whether video models represent and forecast physical laws. We render synthetic videos of simple
mechanical events (collisions, bounces, slides, swings, rolls and tips) in \emph{matched sets}: groups of
videos that share every frame up to an \emph{event} and differ only in what happens after it. The event is the
moment where the physics under test happens: two objects make contact, a ball hits the floor, an object is released,
a block enters a pad or a pendulum enters a damper. In scenes with no such interaction (a projectile in flight, a
pendulum at the top of its swing) it is a chosen frame where the videos branch. \emph{Plausible} videos obey
the law; \emph{implausible} videos break it in a controlled way. A set can hold several plausible outcomes (for
example, collisions that lose different amounts of energy). We pass these videos through video encoders and diffusion transformers and
ask, at several levels, whether the laws hold in what the models compute.

Each law is reduced to dimensionless ratios of one physical quantity measured twice in the same video (for example,
an object's velocity just before and just after a collision). Because both measurements come from the same video and
the same measuring instrument, the ratios do not depend on the video's spatial scale or frame rate
(Section~\ref{sec:laws}). The same ratios are computed from ground-truth labels, from quantities decoded out of model
activations, and from objects tracked in generated videos.

The evaluation levels are:
\begin{itemize}
\item \textbf{L0, state decoding:} are the ingredients of each law (velocity, mass, momentum) linearly decodable from
the model's activations?
\item \textbf{L1, forecasts:} when the model predicts what happens after the event, do the decoded quantities of that
prediction obey the law?
\item \textbf{L2, sensitivity to implausible videos:} is the model's own prediction error larger for the
implausible video than for the plausible one, does it change abruptly at the point where an event stops losing
energy and starts creating it? Does the representation itself change abruptly there, and can a trained classifier
tell implausible from plausible videos?
\item \textbf{L3, invariance:} do the decoded quantities stay unchanged when every object is given the same extra
velocity (Galilean invariance), and when only an object's mass changes in cases where the law says mass should not
matter, such as two balls of different mass falling with the same acceleration (mass independence)?
\item \textbf{L4, steering:} does moving the representation along the direction that encodes a law parameter change
the model's forecast the way the law predicts?
\item \textbf{Generated videos:} do the same ratios hold in videos the diffusion models generate?
\end{itemize}
Every level is reported beside controls that bound what a non-physical shortcut could achieve
(Section~\ref{sec:controls}).

\subsection{Models}

We evaluate two video encoders, V-JEPA~2 ViT-L and VideoMAE-v2-g, and diffusion transformers
(DiTs), including two base/post-trained pairs: CogVideoX-5B vs.\ VideoREPA and Wan2.2-TI2V-5B vs.\ PhysRVG.
V-JEPA~2 sees 64 frames at 24\,fps at $256^2$ (32 temporal tokens of $16\times16$); VideoMAE-v2-g sees 16 frames at
8\,fps at $224^2$; DiTs use their native resolution, frame count and rate. A temporal token is \emph{pre} if all its
frames precede the event, \emph{post} if all follow, and is otherwise excluded. Controls include an untrained V-JEPA~2
and two positive controls trained on plausible renders (Section~\ref{sec:controls}).

\subsection{Videos}
\label{sec:videos}
Videos are rendered with Kubric at 48\,fps, 289 frames (6\,s), $832\times480$, with the event in frames 112--128. Collisions are simulated in PyBullet up to contact; every other scene
is written in closed form. Each implausible video breaks exactly one law. All analyses
of post-event motion use the \emph{analysis window} of a set: from the event to the first later event (an object
leaving the frame, a second contact) in any of its videos. Table~\ref{tab:videos} lists the videos for
each law; some laws are tested in more than one scene.

\begin{table}[h]
\centering\small
\renewcommand{\arraystretch}{1.25}
\begin{tabularx}{\linewidth}{@{}>{\raggedright\arraybackslash}p{1.9cm}>{\raggedright\arraybackslash}p{3.3cm}>{\raggedright\arraybackslash}X>{\raggedright\arraybackslash}X@{}}
\toprule
Law & Scene & Plausible videos & Implausible videos \\
\midrule
Conservation of momentum & two spheres collide on a frictionless plane, seen from above &
collisions at several restitutions; the same collision at five hidden mass ratios $m_A/m_B\in\{0.33,\dots,3\}$ (identical look) &
a drift added to both balls along the contact normal, which changes total momentum but not restitution; in one
comparison each is paired with the plausible video whose post-event kinetic energy is closest, with faster and
slower implausible videos weighted equally, so that kinetic energy alone cannot tell them apart \\
Newton's first law & puck gliding on ice & constant velocity & speed $\times0.7$--$1.3$, or a $20^\circ$ turn, with no
force acting; the main comparison uses one slower and one faster video ($\times0.9$, $\times1.1$) so that speed alone
cannot separate them \\
No creation of energy & sphere collisions; a ball bouncing on the floor; a block crossing a pad between two ramps; a pendulum swinging
through a damper & a \emph{restitution sweep}: videos that differ only in the restitution $\eps$ (the fraction of
approach speed kept after the event), with $\eps\le1$ (energy lost or kept): $e\in\{0,0.1,\dots,1.0\}$ for the
collisions, $\eps\in\{0.50,0.55,\dots,1.00\}$ for the bounce, pad and damper & the same sweep continued past $\eps=1$
(energy created): $e\in\{1.1,\dots,1.5\}$ for the collisions, $\eps\in\{1.05,\dots,1.50\}$ for the bounce, pad and
damper \\
Galilean invariance & the sphere collisions with a common velocity added to both balls, or with the camera moving at the
opposite velocity (on a faintly checkered floor) & collisions with the added velocity or the moving camera, compared
with the same collision without it & none \\
Uniform gravitational acceleration & projectile; dropped balls & falls at two restitutions & gravity $\times1.5$ after the event; horizontal
velocity $\times0.6$; the heavier ball falling faster \\
Coulomb friction & block already sliding on one uniform floor, branching mid-slide; two blocks sliding & constant
deceleration to a stop & sliding back after stopping; deceleration proportional to speed; deceleration
$\times0.7$, $0.85$, $1.15$ or $1.3$; the heavier block decelerating $1.3\times$ faster \\
Rolling without slipping & solid and hollow cylinder released on one incline & the solid cylinder arriving first & the hollow cylinder
arriving first; both arriving together; rolling with slip \\
Torque balance & block resting on a ledge at overhang $x\in\{0.55,\dots,1.45\}$ & tipping iff $x>1$ & tipping when
$x\le1$, or staying when $x>1$ \\
Constant pendulum period & pair of pendulums & unchanged period & period $\times0.8$ or $\times1.25$; the heavier bob swinging
faster \\
Conservation of angular momentum & puck strikes a free bar off-centre & two restitutions & the bar's spin $\times0.8$ or $\times1.2$ with kinetic energy and linear momentum unchanged (the main
comparison uses both together); the bar's spin $\times0.7$--$1.3$ \\
\bottomrule
\end{tabularx}
\caption{Videos by law. All videos in a matched set share every frame before the event.}
\label{tab:videos}
\end{table}

\subsection{Measuring physical laws}
\label{sec:laws}

We test nine laws of classical mechanics:
\begin{itemize}
\item \textbf{Conservation of momentum:} in a two-body collision, the momentum one body loses the other gains; with
no force acting, velocity stays constant (Newton's first law).
\item \textbf{No creation of energy:} a collision, bounce or other exchange cannot create mechanical energy.
\item \textbf{Uniform gravitational acceleration:} gravity is a constant acceleration, independent of horizontal
motion and of mass.
\item \textbf{Coulomb friction:} kinetic friction slows a sliding body at a constant rate, independent of its speed and
mass.
\item \textbf{Rolling without slipping:} on one incline, a body's acceleration depends only on how its mass is
distributed, so a solid cylinder beats a hollow one, and its spin matches its speed.
\item \textbf{Torque balance:} a block on a ledge tips if and only if its centre of mass lies beyond the edge.
\item \textbf{Constant pendulum period:} a pendulum keeps its period, and the period does not depend on the bob's mass.
\item \textbf{Conservation of angular momentum:} a contact between free bodies conserves angular momentum about the
centre of mass.
\item \textbf{Galilean invariance:} adding a common velocity to every body changes neither the restitution nor the
scattering angle in the centre-of-mass frame.
\end{itemize}
To test a law in a video, we compute a \emph{law statistic}: a single number calculated from the objects' measured
motion (positions, velocities, accelerations) that takes a known value whenever the law holds. Each law statistic is a
dimensionless ratio of one quantity measured twice in the same video, for example an object's velocity just before and
just after a collision. Table~\ref{tab:laws} lists the statistics. For each one we give its \emph{physical} value,
the value the law requires, and its \emph{null} value, the value it would take under a simple alternative that ignores
the law (for example, giving both balls in a collision equal and opposite velocity changes regardless of their masses). Scores later place a model between
these two values.

The statistics use the following notation. Each video contains one or two objects, labelled $A$ and $B$, with masses
$m_i$ and radii $r_i$; $k_i=I_i/(m_ir_i^2)$ is an object's inertia factor, which describes how its mass is distributed
($I_i$ is its moment of inertia). Unprimed quantities describe the state just before the event and primed quantities
the state just after it, so $\Delta\mathbf v_i=\mathbf v_i'-\mathbf v_i$ is the change in an object's velocity and
$\mathbf q_i=m_i\Delta\mathbf v_i$ the change in its momentum. In a collision, $\mathbf n$ is the unit vector from
$A$'s centre to $B$'s at contact and $u_n=(\mathbf v_A-\mathbf v_B)\cdot\mathbf n$ is the speed at which they
approach. $M=m_A+m_B$ is the total mass, $\mu=m_Am_B/M$ the reduced mass, $\mathbf V$ the velocity of the centre of
mass, $\ghat$ the direction of gravity and $\eps$ the restitution. Subscripts h and l mark the heavier and lighter
object. Accelerations are obtained by fitting a quadratic, $\mathbf p(t)\approx\mathbf c_0+\mathbf c_1t+\mathbf
c_2t^2$, to the positions over 0.1\,s on each side of the event and taking $\mathbf a=2\mathbf c_2$.

\begin{table}[h]
\centering\small
\renewcommand{\arraystretch}{1.35}
\begin{tabularx}{\linewidth}{@{}l>{\raggedright\arraybackslash}X>{\raggedright\arraybackslash}p{1.7cm}>{\raggedright\arraybackslash}p{3.0cm}@{}}
\toprule
Law & Statistic & Physical & Null \\
\midrule
Cons.\ of momentum & $\eta=\dfrac{\|\mathbf q_A+\mathbf q_B\|}{\|\mathbf q_A\|+\|\mathbf q_B\|}$;\
antiparallel $=-\cos(\Delta\mathbf v_A,\Delta\mathbf v_B)$;\ normal$_i=\dfrac{|\Delta\mathbf v_i\cdot\mathbf n|}{\|\Delta\mathbf v_i\|}$
& $0;\,1;\,1$ & mass-blind: $\eta_{\mathrm{blind}}=\frac{|m_A-m_B|}{m_A+m_B}$ \\
Newton's first law & $\|\mathbf V'\|/\|\mathbf V\|$;\ $\frac1\pi\max_i\angle(\mathbf v_i,\mathbf v_i')$ & $1;\,0$ & stop: $0$ \\
No energy creation & collision $\eps=-\dfrac{(\mathbf v_A'-\mathbf v_B')\cdot\mathbf n}{(\mathbf v_A-\mathbf v_B)\cdot\mathbf n}$;
bounce $-\dfrac{\sum\mathbf v_i'\cdot\ghat}{\sum\mathbf v_i\cdot\ghat}$; pad $\dfrac{\|\mathbf V'\|}{\|\mathbf V\|}$;
swing $\sqrt{h'/h}$ & $\eps\le1$ & pass-through \\
Uniform gravity & $g_{\mathrm{ratio}}=\dfrac{\sum\mathbf a_i'\cdot\ghat}{\sum\mathbf a_i\cdot\ghat}$;
axis $=\dfrac{\|H\mathbf V'\|}{\|H\mathbf V\|}$; mass $=\dfrac{\mathbf a_h'\cdot\ghat}{\mathbf a_l'\cdot\ghat}$ & $1;\,1;\,1$ & no gravity: $0;\,0;\,m_h/m_l$ \\
Coulomb friction & decel $=\dfrac{\sum\mathbf a_i'\cdot\hat{\mathbf v}_i'}{\sum\mathbf a_i\cdot\hat{\mathbf v}_i}$ (after/before the event, while moving);
stop $=\max_i\|\mathbf v_{i,\mathrm{stop}}\|/\|\mathbf v_i\|$; mass $=$ heavy/light deceleration & $1;\,0;\,1$ & Stokes drag \\
Rolling w/o slipping & order $=\dfrac{\log(a_{\mathrm{solid}}/a_{\mathrm{hollow}})}{\log\big((1+k_h)/(1+k_s)\big)}$;
slip $=\dfrac{\|\boldsymbol\omega'\|r}{\|\mathbf v'\|}$ & $1;\,1$ & inertia ignored: $0;\,0$ \\
Torque balance & tip $=\ind[(\varphi>0.05)=(x>1)]$; $\alpha$-ratio $=\|\boldsymbol\alpha'\|/\alpha_{\mathrm{phys}}$,
$\alpha_{\mathrm{phys}}=\frac{g_{\mathrm{clip}}}{a}\frac{d}{d^2+(1+\beta^2)/3}$ & $1;\,1$ & block stays: $\alpha$-ratio $0$ \\
Pendulum period & $\exp(\mathrm{mean}_i\log T_i'/T_i)$;\ $T_h'/T_l'$ & $1;\,1$ & $T\propto m^{-1/2}$ \\
Cons.\ of ang.\ mom. & $L_z'/L_z$, $L_z=\sum m_i(\mathbf p_i-\mathbf R)\times(\mathbf v_i-\mathbf V)+\sum I_i\omega_{z,i}$ & $1$ & spin ignored \\
Galilean invariance & $d_\eps=|\eps-\eps_{\mathrm{ref}}|$;\ $d_\theta=|\mathrm{wrap}(\theta_{\mathrm{cm}}-\theta_{\mathrm{cm,ref}})|/\pi$ (vs.\ unboosted video) & $0;\,0$ & pass-through \\
\bottomrule
\end{tabularx}
\caption{Law statistics, with the value each takes when the law holds (physical) and under the no-physics alternative (null).}
\label{tab:laws}
\end{table}

\noindent For statics, $x$ is the overhang (tips iff $x>1$), $d=|x-1|$, $\beta$ the aspect ratio, $a$ the half-length, $\varphi$ the
rotation over the tip window $T_w=\min\big(2\sqrt{a/g_{\mathrm{clip}}},\ \text{time until the block leaves the ledge}\big)$
after release, and $g_{\mathrm{clip}}$ is measured in the same video (no metric $g$ is used).

Because every law statistic is a ratio of two measurements of the same kind, the units of measurement cancel.
This matters because the quantities are measured in pixels and frames, or decoded by a probe, and none of these
share a common physical scale with the real scene. Proposition~\ref{prop:inv} makes this precise. It considers three
changes: measuring lengths in different units (pixels instead of metres, a factor $s$), playing the video at a
different speed (a factor $\tau$ on time), and a probe that systematically reads every velocity and acceleration too
small or too large by the same factor $\gamma$.

\begin{proposition}[Invariance]
\label{prop:inv}
Every law statistic is unchanged when lengths are rescaled by any $s>0$ or time by any $\tau>0$. When every measured
velocity and acceleration is multiplied by a common factor $\gamma$, every statistic is unchanged except three, which
compare quantities that a probe would read with different instruments: slip and $\alpha$-ratio (each multiplied by
$\gamma^{-1}$) and $L_z'/L_z$ (which changes by an amount that depends on how the angular momentum is split between
orbital motion and spin).
\end{proposition}

\noindent In practice, a probe with a uniform reading error cannot by itself make a model look as if it obeys or
breaks a law, except for those three statistics. Because ridge probes do shrink their predictions, every probe is
calibrated before use (Section~\ref{sec:probe}).



\subsection{Features and forecasts}
\paragraph{Pooling.} $\mathbf z=\sum_{t,u,v}w(t,u,v)\,\mathbf a(t,u,v,:)$ with $w$ a normalised binary mask over a
time window (pre/post/all): \emph{frame} (all tokens), \emph{ROI} (tokens within $1.5$ equivalent radii of the
object), \emph{object pair} (both objects), \emph{contact} (last pre-event ROIs dilated by 2 tokens, fixed in time). ROI
masks use each video's true post-event object position, so every ROI result is reported with a \emph{mask-only
baseline} (the same probe trained on the mask weights alone).

\paragraph{Forecast and prediction error.}
\begin{itemize}
\item \textbf{V-JEPA~2:} the encoder sees the context frames only; the predictor fills the post tokens. Error (L1,
MSE, $1-\cos$) is against the encoder's output on the full video.
\item \textbf{VideoMAE-v2:} future tokens are masked; the decoder predicts normalised pixels. Error is the
pre-training MSE and a de-normalised (raw-pixel) variant.
\item \textbf{DiTs:} error is the training loss on post tokens, averaged over SNR $\in\{-20,-8,0,8\}$\,dB and two
fixed noise draws: flow models $(\hat{\mathbf v}-(\boldsymbol\epsilon-\mathbf z_0))^2$, DDPM
$(\hat{\mathbf z}_0-\mathbf z_0)^2/(1-\bar\alpha_t)$. It is computed with the context left clean and with the whole video noised.
The forecast is a sampled continuation from a fixed seed.
\end{itemize}

\paragraph{Forecast probe.} Encoders: one probe on real pre-event features (encoder space) and a separate probe on the
predictor's outputs for the real post-event frames; the forecast is read by the latter (laws read per timestep use one
probe on real pre and post features of the encoder's final output, the predictor's target space). DiTs: one
probe trained on real pre and post features of the same layer. No probe is ever trained on forecasts.

\subsection{Probes}
\label{sec:probe}
\paragraph{Ridge regression.} Standardised features, centred targets,
$\hat W(\lambda)=V\,\mathrm{diag}\big(\frac{\sigma_j}{\sigma_j^2+\lambda}\big)U^\top(Y-\bar Y)$,
$\lambda\in\{10^{-2},\dots,10^5\}$ chosen by inner grouped 5-fold CV. Folds are split by matched set. Score:
pooled out-of-fold $R^2=1-\sum(y-\hat y)^2/\sum(y-\bar y)^2$, averaged over outputs.

\paragraph{Calibrated probe (R1).} Ridge shrinks predictions toward the mean, which biases ratios. A per-output affine
map $f(\mathbf z)\approx c+d\,y$ is fit on inner out-of-fold predictions and inverted: $\mathrm{decode}(\mathbf z)=(f(\mathbf z)-c)/d$,
applied only where $d$ is significant ($t\ge3$; otherwise the raw prediction is kept). The probe is
cross-fitted (each video decoded by a probe that never saw its set), and, where pre and post features share a space,
\emph{joint}: a video's pre and post rows go through the same map, so both terms of a ratio are read by the same probe.

\paragraph{Selectivity.} Following \citet{hewitt2019control}, $\Delta_{\mathrm{sel}}=R^2_{\mathrm{real}}-
R^2_{\mathrm{control}}$ with two control tasks: \emph{scene} (each set receives another set's targets) and
\emph{object} (each object identity gets one fixed random target). For per-timestep rows, features and targets are
first residualised within (video slot, timestep) strata, leaving out the row's own set.

\subsection{Evaluation levels}

\paragraph{L0: state decoding.}
L0 asks whether the quantities that the law statistics are built from can be read linearly out of a model's
activations. For every object, before and after the event, a ridge probe is trained to predict its velocity, speed,
direction of motion, $\log m$, radius and momentum $m\mathbf v$. We also decode \emph{partial momentum}: the part of
$m\mathbf v$ that no linear combination of $\mathbf v$, $m$ and $\log m$ explains, which tests whether the model
combines mass and velocity instead of encoding them separately. Each law then adds the quantities its own statistics
need: release and apex heights for the pad, the dropped balls and the bounce, the bob's angle and angular velocity for the
pendulums, the cylinders' angular velocity for rolling, the block's acceleration and angular velocity for torque
balance, and the bar's spin for angular momentum. For each target we report the best layer and pooling, chosen on
held-out sets, beside the mask-only baseline, the untrained model, the noised input latents for DiTs, and the probe's
selectivity. For laws with a collision or exchange, we also check whether restitution computed from the decoded
velocities tracks the true restitution; L1 uses this check to decide whether energy can be measured at all.
Mass-independence and rolling results are marked uninformative unless mass (for rolling, the inertia factor) is
decodable before the event, above the untrained model and with positive selectivity.

\paragraph{L1: does the forecast obey the law?}
L1 asks whether the model's own prediction of what happens after the event obeys each law. For every law, the
calibrated probe decodes the quantities that the law's statistics need from the forecast, and the same probe decodes
them from the real post-event frames. The real frames show what the probe can read when the physics is correct, so
every forecast score is reported beside this real-video reference. Each law's quantities are obtained in the way its
statistic requires. For uniform gravity, rolling, Newton's first law and angular momentum, a line is fitted to the
decoded velocity before and after the event, over the frames in which the objects move freely and are in view (at
least three on each side), giving a velocity and an acceleration on each side. For the bounce, the two lines are
extrapolated to the contact. For Coulomb friction, one line is fitted after the event while the block still moves, and
the mean decoded velocity after the stop shows whether the block stays at rest. For the pendulum period, the period
is read from the times at which the decoded horizontal offset of each bob crosses zero. For torque balance, the decoded
spin is integrated over the tip window $T_w$ to give the tipping angle, and its slope gives the angular acceleration. For the
pad and the pendulum damper, restitution is $\sqrt{h'/h}$ from the highest decoded height before and after the pad or damper.

Each statistic $T$ is then scored on the scale from the law's null value to its physical value,
\[
S=1-\frac{|\log T-\log T_{\mathrm{phys}}|}{|\log T_{\mathrm{null}}-\log T_{\mathrm{phys}}|},
\]
so that $S=1$ when the forecast obeys the law and $S=0$ when it matches the no-physics alternative. When the null
value is $0$, the same score is computed without logarithms, and when the null is undefined we use
$1-|T-T_{\mathrm{phys}}|$. Each video's score is clipped to $[-1,1]$, and a video whose decoded statistic is undefined
(or not positive in the logarithmic form) scores $-1$. The reported score is the median over plausible videos.

Two laws need small changes to this recipe. For conservation of momentum, the imbalance $\eta$ is undefined when the
objects exchange no momentum, so only videos in which both the real frames and the forecast show a real exchange are
kept, and $\eta$ is pooled over those videos before scoring (velocity changes use decoded pre- and post-event
velocities, read by the same probe; the hidden-mass sets, whose pre-event frames do not show the mass ratio, are not scored):
\[
S_{\mathrm{mom}}=\mathrm{clip}_{[-1,1]}\!\left(1-\frac{\eta^{\mathrm{fc}}-\eta^{\mathrm{real}}}{\eta^{\mathrm{blind}}-\eta^{\mathrm{real}}}\right),
\]
where $\eta^{\mathrm{blind}}$ is the mass-blind null. For no creation of energy, the law is a one-sided bound
($\eps\le1$), so any forecast that loses energy satisfies it. We therefore report how often the forecast creates more
energy than the probe's own error on real videos can explain, $P_{\mathrm{create}}=\Pr(\hat\eps>1+\tau)$ with $\tau$
the 95th percentile of that error, and we treat energy as measurable only when restitution decoded from real videos
tracks the true restitution ($R^2\ge0.5$). Where a filter depends on the forecast, we also report the share of videos
that pass it and a score that counts filtered videos as failures. A law claim requires the forecast to beat three
negative-control forecasts read by the same probe (the last context frame repeated, another set's forecast, and
another set's real outcome) and, for ROI pools, the same score computed from the mask weights alone.

\paragraph{L2: sensitivity to implausible videos.}
L2 asks whether a model distinguishes implausible from plausible videos. It uses three signals: the model's own
prediction error $\mathcal S$, averaged over the analysis window, which needs no training; the model's activations;
and a trained classifier.

For each law, a matched set gives a comparison between its implausible and plausible videos. Conservation of momentum
compares collisions with and without the momentum drift, at the same restitution. Newton's first law compares the
gliding puck that changes speed or turns with the one that keeps its velocity. Coulomb friction compares blocks that
slide back after stopping, slow in proportion to their speed, decelerate at another rate, or whose heavier block
decelerates faster, with the plausible slide. The pendulum period compares pendulums
whose period changes, or whose heavier bob swings faster, with ones that keep their period. Conservation of angular
momentum compares bars that spin too fast or too slow (in the main comparison with energy unchanged) with the plausible
bar. Uniform gravity compares falls with stronger gravity, reduced horizontal velocity, or a heavier ball falling
faster with the plausible fall. Rolling without slipping compares the hollow cylinder arriving first, a tie, or rolling
with slip with the plausible race. Galilean invariance is tested in L3. For each comparison, $d=\mathcal
S(\text{implausible})-\mathcal S(\text{plausible})$ is computed in every set, and the win rate
$\mathrm{WR}=\frac1N\sum(\ind[d>0]+\frac12\ind[d=0])$ is the fraction of sets in which the implausible video is more
surprising. Each comparison also reports every implausible video alone and flags implausible videos that also create
energy. Because an implausible video may simply move differently, the same win rate is computed for two pixel
baselines (repeating the last frame, and continuing each object at constant velocity), and we report the model's
margin over them, $\Delta\mathrm{WR}=\mathrm{WR}_{\mathrm{model}}-\mathrm{WR}_{\mathrm{pixel}}$. Comparisons between two
plausible videos that differ only in motion serve as controls where a family has two plausible outcomes.

Two laws have a boundary instead of a single implausible video, so they are tested by whether the prediction error
changes abruptly at that boundary. For no creation of energy, the boundary is $\eps=1$ along a restitution sweep. In
each set the log prediction error is standardised to $s_g(\eps)$, its curvature is the second difference
$D^2_g(\eps)=s_g(\eps+h)-2s_g(\eps)+s_g(\eps-h)$, and the boundary score is the curvature at $\eps=1$ above a smooth trend
$p_g$ fitted away from the boundary,
\[
K_g=D^2_g(1)-p_g(1),\qquad \bar K=\mathrm{mean}_g K_g .
\]
A model that encodes the law changes its error abruptly at $\eps=1$, while any purely kinematic signal changes
smoothly there. The trend excludes points near the model's own expected restitution (estimated with the set left out) and near points
where a kinematic signal has a minimum, a stop or a speed match. The result counts as positive when $\bar K$ is above zero, above both pixel baselines and
indistinguishable from zero for the untrained model, and as inconclusive when the model's expected restitution lies
within $0.2$ of the boundary. For torque balance, the boundary is the overhang $x=1$: we compute $c_g=\log\mathcal
S(\text{implausible})-\log\mathcal S(\text{plausible})$ at each overhang, which the law predicts is positive on both
sides of $x=1$, and apply the same curvature test across overhangs.

The same energy boundary test is also applied to the activations themselves. With $\mathbf a(\eps)$ the pooled
post-event activation, the curvature is $N(\eps)=\|\mathbf a(\eps+h)-2\mathbf a(\eps)+\mathbf a(\eps-h)\|$ and the score
is $K^{\mathrm{rep}}_g=\big(N(1)-p(1)\big)/\mathrm{median}\,N$, compared with the same score computed on pixels, on the
true kinematics and on the untrained model.

Finally, a logistic-regression classifier is trained on post-event activations to separate implausible from
plausible videos, and scored by out-of-fold AUC. This applies to the laws whose violations are equalities
(conservation of momentum, Newton's first law, uniform gravity, Coulomb friction, rolling without slipping, the
pendulum period and conservation of angular momentum); energy and torque balance are covered by the boundary tests above. Each AUC is compared with the same
classifier trained on the true post-event kinematics (speeds, velocity changes and kinetic energy), which shows how
much the model separates beyond what the motion alone allows; a positive result must beat it and, for object-patch
pools, the mask-only baseline.

\paragraph{L3: invariance.}
L3 asks whether the quantities a model represents respect the two invariance laws. For Galilean invariance, velocity
probes trained on collisions at every added velocity decode each collision's restitution and centre-of-mass scattering
angle from the forecast, and we measure how much these change ($d_\eps$, $d_\theta$) relative to the forecast of the same
collision without the added velocity, and between a moving camera and the same velocity added to the balls; the same
measurement on real videos is the floor. An added-velocity level is reported as unmeasurable when the probe's restitution error on real videos
exceeds $0.25$. Mass independence is tested for three laws whose outcomes should not depend on mass: uniform gravity
(dropped balls), Coulomb friction (two sliding cubes) and the pendulum period (pendulum pair). For each, one calibrated probe shared by both objects decodes
the heavy-to-light ratio $t$ of the relevant quantity (acceleration, deceleration or period) from the forecast and from
the real frames, and the score
\[
S=1-\frac{|\log t_{\mathrm{fc}}-\log t_{\mathrm{real}}|}{|\log t_{\mathrm{implaus}}-\log t_{\mathrm{real}}|}
\]
places the forecast between the real plausible videos and the real implausible ones; it must also beat the same score
computed from the ROI mask weights alone. We also test whether probes
transfer from clean to cluttered renders, using object size as a control target.

\paragraph{L4: steering.}
L4 asks whether a model uses what it represents. For each law we choose a parameter $\theta$ that the law ties to an
observable outcome (Table~\ref{tab:steer}), find the direction in activation space along which $\theta$ is encoded,
move the activations along it, rerun the model's forecast, and check whether the decoded outcome changes as the law
predicts. The direction is a curve $\mathbf c(\theta)$ fitted to the activations of real videos; steering by $\delta$
adds $\mathbf c(\theta_{\mathrm{ref}}+\delta)-\mathbf c(\theta_{\mathrm{ref}})$ to the activations. The curve is fitted
where the shift is added: on the context tokens of an encoder layer, or, for energy, on the post-event tokens of an
intermediate predictor or decoder block. For DiTs it is fitted separately at each noise level, and each sampling step
uses the curve of the nearest noise level (or no shift if that curve failed). Steering is run only where the curve
passes a precondition: it must predict held-out scenes (leave-one-scene-out $R^2$) and, where $\theta$ varies within a
scene, each scene's offset along it must survive averaging over scenes, both better than with $\theta$ shuffled
within scenes (across scenes where each scene has one $\theta$).

\begin{table}[h]
\centering\small
\begin{tabularx}{\linewidth}{@{}l>{\raggedright\arraybackslash}p{3.3cm}>{\raggedright\arraybackslash}p{3.6cm}>{\raggedright\arraybackslash}X@{}}
\toprule
Law & Steered parameter $\theta$ & Readout & Law's prediction \\
\midrule
Cons.\ of momentum & $\log(m_A/m_B)$ & log ratio of $B$'s to $A$'s velocity change along $\mathbf n$, fitted with errors in both (contact region) & slope 1 \\
No energy creation & restitution $\eps$ & decoded $\hat\eps$ & rises, then saturates at 1 \\
Uniform gravity & $\log m_A$ & log heavy/light acceleration ratio & slope 0 \\
Coulomb friction & $\log m_A$ & $\log$ of $A$'s deceleration & slope 0 \\
Rolling w/o slipping & inertia factor $k_A$ & $\log$ of $A$'s acceleration & slope $-1/(1+k)$ \\
Torque balance & overhang $x$ & whether the block tips & step at $x=1$ \\
Pendulum period & $\log m_A$ & $\log$ of $A$'s period & slope 0 \\
Cons.\ of ang.\ mom. & impact point on the bar & $L_z'/L_z$ & slope 0 \\
Galilean invariance & added velocity & decoded $\hat\eps$ & slope 0 \\
\bottomrule
\end{tabularx}
\caption{Steering tests. Laws with a predicted slope of $0$ are informative only when the push moves the model's
readout of $\theta$ or a control direction does move the readout; otherwise the push is reported as inert.}
\label{tab:steer}
\end{table}

We also run an interchange intervention: the component of one video's activations along the curve is replaced by
that of another video with a different $\theta$. For laws that predict a change, we measure how often the readout moves
closer to the law's predicted change than to no change (IIA) and the fraction of the required change it achieves. For
laws that predict no change, we measure the change of the outcome relative to a reference (how much the readout
differs between real videos that the law says do differ, or else its spread over real videos), and check that the
model's own readout of $\theta$ did move (for DiTs, read from the activations at the swapped layer), so that an unchanged
outcome is not simply a failed swap. For momentum we also
patch the collision directly: object $A$'s pre-contact tokens are replaced by those of a donor video with the same
contact normal and contact time step (for DiTs, the donor runs through the same denoising step and noise in a paired
pass), and the change in $B$'s outcome is scored by the error ratio $|\text{required}-\text{observed}|/|\text{required}|$
(1 for no change) with a sign-flip test. The ratio of the response slopes in high- and low-impulse pairs separates
physics, which predicts their impulse ratio, from a constant leak, which predicts 1; the patch passes only if the
change is significantly closer to the required one than to none, beats the controls, and this ratio excludes 1.

Every steering result is compared with controls (passive propagation: the shifted activations read at the steered
block without the model's downstream computation; one random direction of the same size; and a direction encoding
object speed for energy or object size for the mass laws) and with the change produced by showing the same
$\Delta\theta$ through a real video. We also report the per-scene effects, the share of scenes that move the wrong way,
and how far the steered activations lie from real ones.

For models that pass steering, the interchange test or the L2 boundary test, we localise the law: the swap IIA is
computed by layer (and noise level), by token group (each object, the contact region, the background, the frames
before or after the event) and by component (attention heads and MLPs: an attribution-patching pre-screen, exact
patching, and knock-out by mean ablation; for energy, the knock-out also tests whether removing a component lets the
forecast create energy), and a subspace learned by distributed alignment search is compared with the fitted curve. A
layer or component counts as a location only when both the swap and the knock-out beat their controls (for the swap:
a random direction, shuffled targets, a size direction, passive propagation and the untrained model; for the
knock-out: random directions or components of the same size).

\paragraph{Generated videos.}
For the diffusion models, we also test the laws in the videos they generate. A colour tracker, with a marker-angle
tracker for the rolling cylinders and the struck bar, follows each object, and the same law statistics are computed
from the tracks without using any metric scale or frame rate. Each tracker is first validated on real plausible and
implausible videos. Generations either continue the real frames up to the event or start from the first frame only;
the first are compared with the real plausible video, and the second are scored for consistency with themselves.

\subsection{Controls}
\label{sec:controls}
Untrained V-JEPA~2; positive controls (a small predictor and a fine-tuned DiT trained on plausible renders); real-video
reference; mask-only baseline; selectivity; mass-blind rule; pixel baselines; speed-matched implausible videos;
plausible-vs-plausible pairs; negative-control forecasts (L1); kinematic baseline (L2 trained detection); steering controls (L4).

\end{document}
